\documentclass[12pt,a4paper]{amsart}

\usepackage{amsmath,amssymb,amsthm}
\usepackage{mathtools}
\usepackage{enumerate}
\usepackage{hyperref}
\usepackage{cleveref}
\usepackage{geometry}
\geometry{margin=2.5cm}

%%------------------------------------------------------------------
%% Theorem environments
%%------------------------------------------------------------------
\newtheorem{theorem}{Theorem}[section]
\newtheorem{lemma}[theorem]{Lemma}
\newtheorem{corollary}[theorem]{Corollary}
\newtheorem{proposition}[theorem]{Proposition}
\newtheorem{definition}[theorem]{Definition}
\newtheorem{remark}[theorem]{Remark}
\newtheorem{assumption}[theorem]{Assumption}

%%------------------------------------------------------------------
%% Notation macros
%%------------------------------------------------------------------
\newcommand{\psin}{\psi_n}
\newcommand{\wncl}{\omega_n^{\mathrm{cl}}}
\newcommand{\rhocl}{\rho^{\mathrm{cl}}}
\newcommand{\IB}{I_\delta}

\begin{document}

\title[Semiclassical Equidistribution of PSWF Densities]%
{Semiclassical Equidistribution of Prolate Spheroidal Wave Function
Densities via Pr\"ufer Phase Analysis}

\author{Anonymous}
\date{April 2026}
\maketitle

%%=================================================================
\begin{abstract}
We prove that the squared densities $\psin^2$ of the prolate
spheroidal wave functions (PSWFs) equidistribute against the
classical density $\rhocl$ on $[-T,T]$ at a rate $O(1/n)$
for all $n \le \gamma N$, where $N = cT/\pi$ is the Shannon number.
The proof uses Pr\"ufer phase analysis, integration by parts for
the oscillatory terms, and an explicit drift cancellation identity
that reduces the amplitude deviation to a single oscillatory
integral. No microlocal machinery, Airy-function matching, or
Riemann--Hilbert methods are used.
This result supplies the semiclassical equidistribution hypothesis
required by the bulk tail bound program of \cite{PaperIII}.
\end{abstract}

%%=================================================================
\section{Introduction}\label{sec:intro}
%%=================================================================

\subsection{Background and motivation}

Prolate spheroidal wave functions $\psin$ are the eigenfunctions
of the finite Fourier operator $\mathcal{F}_T$ on $[-T,T]$:
\[
  \int_{-T}^T \psin(y)\,e^{icxy}\,dy = \mu_n(c)\,\psin(x),
\]
where $c > 0$ is the bandwidth parameter and
$\lambda_n(c) = |\mu_n(c)|^2$ are the eigenvalues.
Equivalently, $\psin$ are eigenfunctions of the differential
operator
\[
  \mathcal{L}_c\psin
  = \bigl[(1-x^2/T^2)\psin'\bigr]'
    + \bigl[c^2(1-x^2/T^2) - \chi_n(c)/T^2\bigr]\psin = 0,
\]
with eigenvalues $\chi_n(c)$ satisfying
$\chi_n(c) \sim n(n+1) + (2n+1)c/\pi + O(c^2/n)$
for $n$ in the bulk regime $n \le \gamma N$, $N = cT/\pi$.

The present paper establishes a quantitative semiclassical
equidistribution theorem for $\psin^2$, which is a key analytic
ingredient in the bulk tail bound program developed in
\cite{PaperI,PaperII,PaperIII}.

\subsection{Main result (informal)}

For $n$ in the bulk regime $n \le \gamma N$, the squared density
$\psin^2$, when renormalized by $\lambda_n(c)$, converges weakly
to the classical equilibrium density
\[
  \rhocl(x) =
  \frac{1}{\pi\sqrt{x_+(n)^2 - x^2}}
  \cdot \frac{1}{\int_{-x_+(n)}^{x_+(n)}
                  \frac{dt}{\sqrt{x_+(n)^2-t^2}}},
  \quad |x| < x_+(n),
\]
at rate $O(1/n)$, uniformly over $C^1$ test functions and
$n \le \gamma N$.

\subsection{Strategy of proof}

The proof has four layers, executed in Sections
\ref{sec:prufer}--\ref{sec:weaklimit}:

\begin{enumerate}[(i)]
  \item \textbf{Pr\"ufer transform} (Section~\ref{sec:prufer}):
  reduce $\psin^2$ to amplitude $r_n^2$ and phase $\theta_n$.
  \item \textbf{Oscillation control}
  (Section~\ref{sec:oscillation}):
  show $\int h\cos(2\theta_n) = O(\|h\|_{C^1}/n)$ by IBP,
  using the monotonicity $\theta_n' \ge cn > 0$ on $I_\delta$.
  \item \textbf{Amplitude drift} (Section~\ref{sec:amplitude}):
  show $r_n^2/r_n^{\mathrm{WKB},2} = 1 + O(1/n)$ via the drift
  cancellation identity, with bootstrap closure for the
  self-referential term.
  \item \textbf{Main theorem} (Section~\ref{sec:weaklimit}):
  assemble the above to obtain the rate-$O(1/n)$
  equidistribution.
\end{enumerate}

\subsection{Logical structure and non-circularity}

\begin{center}
\begin{tabular}{lll}
\hline
Step & Uses & Does \emph{not} use \\
\hline
Lemma~\ref{lem:prufer-oscillation-control} &
  $\theta_n' \ge c_{\delta,\gamma}n$ &
  $r_n$, $\Delta_n$ \\
Lemma~\ref{lem:deviation-ode} &
  Pr\"ufer ODE &
  Lemma~\ref{lem:prufer-oscillation-control} \\
Lemma~\ref{lem:amplitude-drift} &
  Lemmas~\ref{lem:prufer-oscillation-control},
  \ref{lem:deviation-ode} &
  Theorem~\ref{thm:weak-limit} \\
Lemma~\ref{lem:normalization} &
  Lemma~\ref{lem:amplitude-drift} &
  Theorem~\ref{thm:weak-limit} \\
Theorem~\ref{thm:weak-limit} &
  Lemmas~\ref{lem:prufer-oscillation-control},
  \ref{lem:amplitude-drift},
  \ref{lem:normalization} &
  --- \\
\hline
\end{tabular}
\end{center}

\noindent
In particular, no step uses Theorem~\ref{thm:weak-limit} as a
hypothesis; the dependency graph is a strict DAG with no cycles.

\subsection{Relation to prior work}

Oscillatory estimates for PSWF densities appear implicitly in
\cite{OsipovRokhlinXiao}, but without quantitative rates.
The Pr\"ufer-based approach is classical for Sturm--Liouville
theory \cite{Levitan}; its adaptation to the PSWF setting
with explicit $O(1/n)$ rates is new.
The result is a PSWF analogue of the Szeg\H{o} equidistribution
theorem for orthogonal polynomial measures \cite{Szego},
proved here by elementary ODE methods rather than
potential-theoretic or Riemann--Hilbert techniques.

%%=================================================================
\section{Setup, Notation, and Preliminary Lemmas}\label{sec:setup}
%%=================================================================

\subsection{Parameters and bulk regime}

Throughout, $c, T > 0$ are fixed, $N := cT/\pi$ is the Shannon
number, and $0 < \gamma < 1$ is a fixed bulk parameter.
We write $n \le \gamma N$ for the bulk regime.
All constants $C_{\delta,\gamma}$, $c_{\delta,\gamma}$ depend
only on $\delta$, $\gamma$, $c$, $T$.

\subsection{Turning points and classical frequency}

For $n \le \gamma N$, the classical turning point is:
\begin{equation}\label{eq:turning-point}
  x_+(n) := \frac{T}{c}\sqrt{\chi_n(c)},
\end{equation}
with $x_+(n) \le (1-\eta_\gamma)T$ for some $\eta_\gamma > 0$
depending only on $\gamma$.

The classical frequency function is:
\begin{equation}\label{eq:wncl-def}
  \wncl(x) := \sqrt{\chi_n(c)/T^2 - c^2 x^2/T^2 + c^2}
  = \sqrt{c^2(1 - x^2/T^2) + \chi_n(c)/T^2},
\end{equation}
which satisfies $\wncl(x) \asymp n$ uniformly on $I_\delta$
(see Lemma~\ref{lem:G-regularity} below).

\subsection{Bulk interval}

Fix $0 < \delta < \eta_\gamma$. Define the bulk interval:
\begin{equation}\label{eq:Idelta}
  I_\delta := [-(1-\delta)T,\,(1-\delta)T].
\end{equation}
Then $[-x_+(n), x_+(n)] \subset I_\delta \subset [-T,T]$
for all $n \le \gamma N$.

\subsection{Classical density}

\begin{definition}[Classical density]\label{def:rho-cl}
\begin{equation}\label{eq:rho-cl}
  \rhocl(x) :=
  \frac{1/\wncl(x)}{\int_{-x_+(n)}^{x_+(n)} dt/\wncl(t)},
  \quad x \in [-x_+(n), x_+(n)].
\end{equation}
\end{definition}

\subsection{Bohr--Sommerfeld quantization}

\begin{lemma}[Bohr--Sommerfeld for PSWFs]%
\label{lem:bohr-sommerfeld-pswf}
For $n \le \gamma N$:
\begin{equation}\label{eq:BS}
  \int_{-x_+(n)}^{x_+(n)} \wncl(x)\,dx
  = \pi\bigl(n + \tfrac{1}{2}\bigr)
    + O\!\left(\frac{1}{n}\right).
\end{equation}
\end{lemma}

\begin{proof}
This is the standard WKB quantization condition for the
Sturm--Liouville eigenvalue problem \cite{Levitan},
specialized to the PSWF differential operator.
The error $O(1/n)$ follows from a uniform WKB estimate,
valid on $I_\delta$ where $\wncl(x) \ge c_{\delta,\gamma}n > 0$;
see \cite{OsipovRokhlinXiao}, Ch.~4.
\end{proof}

\subsection{Regularity of the WKB coefficient}

\begin{lemma}[Regularity of $G$ and $\wncl$]%
\label{lem:G-regularity}
Define:
\begin{equation}\label{eq:G-def}
  G(x) := \frac{(\wncl)'(x)}{2\wncl(x)} + \frac{p'(x)}{4p(x)},
  \qquad p(x) := 1 - x^2/T^2.
\end{equation}
Then on $I_\delta$, uniformly in $n \le \gamma N$:
\begin{enumerate}[(i)]
  \item $\wncl(x) \ge c_{\delta,\gamma}\,n$,
  \item $|(\wncl)'(x)| \le C_{\delta,\gamma}$,
  \item $\|G\|_{L^\infty(I_\delta)} \le C_{\delta,\gamma}/n$,
  \item $\|G'\|_{L^1(I_\delta)} \le C_{\delta,\gamma}/n$,
  \item $p(x) \ge \delta^2$ and
        $|p'(x)| = 2|x|/T^2 \le 2/T$.
\end{enumerate}
\end{lemma}

\begin{proof}
(i): $\wncl(x)^2 = c^2 p(x) + \chi_n/T^2 \ge \chi_n/T^2
\ge C_\gamma n^2/T^2$, using $\chi_n \ge C_\gamma n^2$
for $n \le \gamma N$.
(ii): Direct differentiation of \eqref{eq:wncl-def}.
(iii)--(iv): $G = (\wncl)'/(2\wncl) + p'/(4p)$;
each term is $O(1/n)$ on $I_\delta$ by (i), (ii), (v).
(v): Explicit from $p(x) = 1 - x^2/T^2 \ge \delta^2$
on $I_\delta$.
\end{proof}

\subsection{Turning point cutoff}

\begin{lemma}[Exponential turning point cutoff]%
\label{lem:turning-point-cutoff}
For $n \le \gamma N$ and $f \in L^\infty([-T,T])$:
\begin{equation}\label{eq:cutoff}
  \left|\int_{[-T,T]\setminus I_\delta}
    f(x)\,\psin(x)^2\,dx\right|
  \le C_{\delta,\gamma}\,e^{-\beta_{\delta,\gamma} n}\,
     \lambda_n\,\|f\|_{L^\infty},
\end{equation}
where $\beta_{\delta,\gamma} > 0$ depends only on
$\delta$, $\gamma$, $c$, $T$.
\end{lemma}

\begin{proof}
Standard: $\psin$ decays exponentially outside the
classically allowed region $[-x_+(n), x_+(n)]$,
and since $x_+(n) \le (1-\eta_\gamma)T < (1-\delta)T$,
the region $[-T,T]\setminus I_\delta$ lies strictly outside
the turning points; see \cite{OsipovRokhlinXiao}, Ch.~5.
\end{proof}

%%=================================================================
\section{Pr\"ufer Transform}\label{sec:prufer}
%%=================================================================

\subsection{Pr\"ufer substitution}

On $I_\delta$, $\psin$ satisfies:
\begin{equation}\label{eq:SL}
  (p(x)\psin')' + \wncl(x)^2/p(x)\cdot\psin = 0.
\end{equation}
Define the Pr\"ufer variables $(r_n, \theta_n)$ by:
\begin{equation}\label{eq:prufer-def}
  \psin(x) = r_n(x)\sin(\theta_n(x)),
  \quad
  p(x)\psin'(x) = r_n(x)\wncl(x)\cos(\theta_n(x)).
\end{equation}

\subsection{Pr\"ufer ODEs}

\begin{lemma}[Pr\"ufer system]\label{lem:prufer-ode}
The Pr\"ufer variables satisfy on $I_\delta$:
\begin{align}
  \theta_n'
  &= \wncl(x)/p(x) + G(x)\sin(2\theta_n(x)),
  \label{eq:theta-ode}\\
  (\log r_n)'
  &= G(x)\cos(2\theta_n(x)) + D(x),
  \label{eq:r-ode}
\end{align}
where $D(x) := p'(x)/(4p(x))$ and $G(x)$ is as in
Lemma~\ref{lem:G-regularity}.
\end{lemma}

\begin{proof}
Standard Pr\"ufer differentiation; see \cite{Levitan}, Ch.~1.
\end{proof}

\subsection{Phase monotonicity}

\begin{lemma}[Phase lower bound]\label{lem:phase-lower}
On $I_\delta$, uniformly in $n \le \gamma N$:
\begin{equation}\label{eq:phase-lb}
  \theta_n'(x) \ge \frac{c_{\delta,\gamma}\,n}{T}.
\end{equation}
\end{lemma}

\begin{proof}
From \eqref{eq:theta-ode} and
Lemma~\ref{lem:G-regularity}(i),(iii):
$\theta_n' \ge \wncl/p - |G|
\ge c_{\delta,\gamma}n - C_{\delta,\gamma}/n
\ge c_{\delta,\gamma}n/2$
for $n$ sufficiently large;
adjusting the constant absorbs the finite cases.
\end{proof}

\subsection{Phase second derivative bound}

\begin{lemma}[Phase second derivative]%
\label{lem:phase-second-deriv}
On $I_\delta$, uniformly in $n \le \gamma N$:
\begin{equation}\label{eq:phase-2nd}
  |\theta_n''(x)| \le C_{\delta,\gamma},
  \qquad
  \left|\frac{\theta_n''(x)}{(\theta_n'(x))^2}\right|
  \le \frac{C_{\delta,\gamma}}{n^2}.
\end{equation}
\end{lemma}

\begin{proof}
Differentiate \eqref{eq:theta-ode}:
$\theta_n'' = (\wncl/p)' + (G\sin(2\theta_n))'$.
Using Lemma~\ref{lem:G-regularity},
$(\wncl/p)' = O(1)$,
$G' = O(1/n)$,
$G\theta_n' = O(1)$,
so $|\theta_n''| \le C_{\delta,\gamma}$.
The ratio bound follows from
$\theta_n' \ge c_{\delta,\gamma}n$
(Lemma~\ref{lem:phase-lower}).
\end{proof}

%%=================================================================
\section{Oscillation Control Lemma}\label{sec:oscillation}
%%=================================================================

\begin{lemma}[Pr\"ufer oscillation control, uniform in subintervals]%
\label{lem:prufer-oscillation-control}
Let $h \in C^1(I_\delta)$. Then for every $x_0, x \in I_\delta$:
\begin{equation}\label{eq:osc-control-sup}
  \sup_{x_0, x \in I_\delta}
  \left|\int_{x_0}^{x} h(t)\cos(2\theta_n(t))\,dt\right|
  \le \frac{C_{\delta,\gamma}}{n}
      \left(\|h\|_{L^\infty(I_\delta)}
        + \|h'\|_{L^1(I_\delta)}\right).
\end{equation}
In particular:
\begin{equation}\label{eq:osc-control}
  \left|\int_{I_\delta} h(x)\cos(2\theta_n(x))\,dx\right|
  \le \frac{C_{\delta,\gamma}}{n}
      \left(\|h\|_{L^\infty(I_\delta)}
        + \|h'\|_{L^1(I_\delta)}\right).
\end{equation}
\end{lemma}

\begin{proof}
Fix $x_0, x \in I_\delta$ with $x_0 < x$.
Since $\theta_n' \ge (c_{\delta,\gamma}/T)n > 0$ on $I_\delta$
by Lemma~\ref{lem:phase-lower},
integrate by parts:
\begin{equation}\label{eq:ibp2}
  \int_{x_0}^{x} h\cos(2\theta_n)\,dt
  = \left[\frac{h(t)\sin(2\theta_n(t))}{2\theta_n'(t)}\right]_{x_0}^{x}
    - \int_{x_0}^{x} \sin(2\theta_n)\cdot
      \frac{d}{dt}\!\left(\frac{h}{2\theta_n'}\right)\,dt.
\end{equation}

\noindent\textit{Boundary term.}
$|\sin(2\theta_n)| \le 1$ and
$\theta_n' \ge (c_{\delta,\gamma}/T)n$,
so the boundary term is bounded by
$T\|h\|_{L^\infty}/(c_{\delta,\gamma}n)$,
uniformly in $x_0, x \in I_\delta$.

\noindent\textit{Interior: $h'$-term.}
\[
  \left|\int_{x_0}^x \sin(2\theta_n)\frac{h'}{2\theta_n'}\,dt\right|
  \le \frac{T}{2c_{\delta,\gamma}n}\|h'\|_{L^1(I_\delta)}.
\]

\noindent\textit{Interior: $\theta_n''$-term.}
By Lemma~\ref{lem:phase-second-deriv},
$|\theta_n''(t)/(\theta_n'(t))^2| \le C_{\delta,\gamma}/n^2$,
so:
\[
  \left|\int_{x_0}^x\sin(2\theta_n)
    \frac{h(t)\theta_n''(t)}{2(\theta_n'(t))^2}\,dt\right|
  \le \frac{C_{\delta,\gamma}\,|I_\delta|}{n^2}
      \|h\|_{L^\infty(I_\delta)}
  \le \frac{2TC_{\delta,\gamma}}{n^2}
      \|h\|_{L^\infty(I_\delta)}.
\]

\noindent\textit{Assembly.}
All three bounds are uniform in $x_0, x \in I_\delta$,
giving \eqref{eq:osc-control-sup}.
The bound \eqref{eq:osc-control} follows by taking
$x_0 = -(1-\delta)T$ and $x = (1-\delta)T$.
\end{proof}

\begin{remark}[Uniformity in $x_0, x$ is essential]
The sup-form \eqref{eq:osc-control-sup} is required in
Lemma~\ref{lem:amplitude-drift}, Step~3, where one needs
$\sup_x|\int_{x_0}^x G\cos(2\theta_n)| \le C/n$
for the exponentiation step to be valid uniformly in $x \in I_\delta$.
A bound only for the full integral over $I_\delta$ would be
insufficient here.
\end{remark}

\begin{remark}[Sharpness]
For $h \equiv 1$ and $\theta_n(x) \approx nx/T$,
direct computation gives
$|\int_{x_0}^x\cos(2nt/T)\,dt|
= |{\sin(2nx/T) - \sin(2nx_0/T)}|/(2n/T) = O(T/n)$,
so the bound \eqref{eq:osc-control-sup} is sharp in $n$.
\end{remark}

\begin{corollary}[Oscillation control with slowly varying weight]%
\label{cor:osc-slowly-varying}
If $a \in C^1(I_\delta)$ and $b_n \in C^1(I_\delta)$ satisfies
$\|b_n\|_{L^\infty(I_\delta)} \le C$
and $\|b_n'\|_{L^\infty(I_\delta)} \le C/n$
for a constant $C$ independent of $n$, then:
\begin{equation}\label{eq:osc-sv}
  \left|\int_{I_\delta}
    a(x)\,b_n(x)\,\cos(2\theta_n(x))\,dx\right|
  = O\!\left(\frac{\|a\|_{C^1(I_\delta)}\,
                   \|b_n\|_{L^\infty(I_\delta)}}{n}\right).
\end{equation}
\end{corollary}

\begin{proof}
Set $h := ab_n \in C^1(I_\delta)$.
Then:
\[
  \|h\|_{L^\infty} \le \|a\|_{L^\infty}\|b_n\|_{L^\infty}
\]
and:
\[
  \|h'\|_{L^1}
  \le |I_\delta|\bigl(\|a'\|_{L^\infty}\|b_n\|_{L^\infty}
      + \|a\|_{L^\infty}\|b_n'\|_{L^\infty}\bigr)
  \le 2T\bigl(\|a'\|_{L^\infty}\|b_n\|_{L^\infty}
      + \|a\|_{L^\infty}\cdot C/n\bigr)
  = O\bigl(\|a\|_{C^1}\|b_n\|_{L^\infty}\bigr).
\]
Apply Lemma~\ref{lem:prufer-oscillation-control}
to $h = ab_n$.
\end{proof}

%%=================================================================
\section{Amplitude Drift Control}\label{sec:amplitude}
%%=================================================================

\begin{definition}[WKB amplitude]\label{def:wkb-amp}
Fix $x_0 \in I_\delta$. Define:
\begin{equation}\label{eq:wkb-amp}
  r_n^{\mathrm{WKB}}(x)
  := r_n(x_0)\sqrt{\frac{\wncl(x_0)}{\wncl(x)}},
  \quad x \in I_\delta.
\end{equation}
This is the unique solution on $I_\delta$ of the averaged
Pr\"ufer amplitude ODE
$(\log r)' = D(x) - (\wncl)'/(2\wncl)$
with initial value $r(x_0) = r_n(x_0)$,
i.e., the Pr\"ufer amplitude ODE \eqref{eq:r-ode}
with the oscillatory $G(x)\cos(2\theta_n)$ term set to zero.
Note that:
\begin{equation}\label{eq:wkb-amp-deriv}
  \frac{d}{dx}\bigl[r_n^{\mathrm{WKB}}(x)^2\bigr]
  = -r_n(x_0)^2\,
    \frac{\wncl(x_0)\cdot(\wncl)'(x)}{\wncl(x)^2}
  = O\!\left(\frac{r_n(x_0)^2}{n}\right),
\end{equation}
uniformly on $I_\delta$ and in $n \le \gamma N$,
since $(\wncl)'(x) = O(1)$ and $\wncl(x) \asymp n$
by Lemma~\ref{lem:G-regularity}(i),(ii).
\end{definition}

\begin{lemma}[Exact amplitude deviation ODE]
\label{lem:deviation-ode}
Let $\Delta_n(x) := r_n(x)^2 - r_n^{\mathrm{WKB}}(x)^2$.
Then $\Delta_n$ satisfies:
\begin{equation}\label{eq:deviation-ode}
  \Delta_n'(x)
  = 2D(x)\,\Delta_n(x)
    + 2r_n^{\mathrm{WKB}}(x)^2\,\widetilde{B}_n(x)
    + 2\Delta_n(x)\,G(x)\cos(2\theta_n(x)),
\end{equation}
where:
\begin{equation}\label{eq:Btilde-def}
  \widetilde{B}_n(x)
  := \frac{(\wncl)'(x)}{2\wncl(x)}\cos(2\theta_n(x))
     + G(x)\cos(2\theta_n(x)).
\end{equation}
In particular $|\widetilde{B}_n(x)| \le C_{\delta,\gamma}/n$
uniformly on $I_\delta$
(Lemma~\ref{lem:G-regularity}(i),(ii),(iii)).
All $p'/(4p)$-contributions cancel exactly as a consequence
of the definition of $r_n^{\mathrm{WKB}}$,
leaving only the oscillatory forcing $\widetilde{B}_n$
with coefficient $O(1/n)$ and the self-referential term
$2\Delta_n G\cos(2\theta_n)$,
which is controlled by the bootstrap argument in
Lemma~\ref{lem:amplitude-drift}, Step~4.
\end{lemma}

\begin{proof}
From \eqref{eq:r-ode}:
$(r_n^2)' = 2r_n^2\bigl(D(x) + G(x)\cos(2\theta_n)\bigr)$.
From Definition~\ref{def:wkb-amp}:
$(r_n^{\mathrm{WKB},2})' = 2r_n^{\mathrm{WKB},2}
\bigl(D(x) - (\wncl)'/(2\wncl)\bigr)$.
Subtracting:
\begin{align*}
  \Delta_n'
  &= 2r_n^2\bigl(D + G\cos(2\theta_n)\bigr)
     - 2r_n^{\mathrm{WKB},2}\bigl(D - (\wncl)'/(2\wncl)\bigr) \\
  &= 2D\bigl(r_n^2 - r_n^{\mathrm{WKB},2}\bigr)
     + 2r_n^2 G\cos(2\theta_n)
     + 2r_n^{\mathrm{WKB},2}(\wncl)'/(2\wncl) \\
  &= 2D\,\Delta_n
     + 2(\Delta_n + r_n^{\mathrm{WKB},2})G\cos(2\theta_n)
     + 2r_n^{\mathrm{WKB},2}(\wncl)'/(2\wncl).
\end{align*}
Collecting the $r_n^{\mathrm{WKB},2}$ terms:
\[
  2r_n^{\mathrm{WKB},2}\bigl[G\cos(2\theta_n)
  + (\wncl)'/(2\wncl)\bigr]
  = 2r_n^{\mathrm{WKB},2}\,\widetilde{B}_n,
\]
with $\widetilde{B}_n$ as in \eqref{eq:Btilde-def},
and the remaining term $2\Delta_n G\cos(2\theta_n)$
is the self-referential contribution
controlled in Lemma~\ref{lem:amplitude-drift}, Step~4.
The $D = p'/(4p)$ terms from both sides cancel exactly,
leaving \eqref{eq:deviation-ode}.
\end{proof}

\begin{lemma}[Amplitude drift]\label{lem:amplitude-drift}
For $n \le \gamma N$ and all $x \in I_\delta$:
\begin{align}
  \left|\frac{r_n(x)^2}{r_n^{\mathrm{WKB}}(x)^2} - 1\right|
  &\le \frac{C_{\delta,\gamma}}{n},
  \label{eq:amp-ratio}\\
  \left|\Delta_n'(x)\right|
  &\le \frac{C_{\delta,\gamma}\,r_n(x_0)^2}{n}.
  \label{eq:amp-dev-deriv}
\end{align}
\end{lemma}

\begin{remark}[Two bounds, two roles]\label{rem:two-bounds}
Bound~\eqref{eq:amp-ratio} gives pointwise ratio control,
used in Lemmas~\ref{lem:normalization}
and~\ref{thm:weak-limit}.
Bound~\eqref{eq:amp-dev-deriv} shows that
$b_n := r_n^2/r_n(x_0)^2$ satisfies
$\|b_n'\|_{L^\infty(I_\delta)} = O(1/n)$,
as required by Corollary~\ref{cor:osc-slowly-varying}.
To see this: $(r_n^2)' = (r_n^{\mathrm{WKB},2})' + \Delta_n'
= O(r_n(x_0)^2/n) + O(r_n(x_0)^2/n)
= O(r_n(x_0)^2/n)$,
using \eqref{eq:wkb-amp-deriv} and \eqref{eq:amp-dev-deriv}.
Dividing by $r_n(x_0)^2$ gives
$\|b_n'\|_{L^\infty} = O(1/n)$.
\end{remark}

\begin{proof}
\textit{Step 1: Drift cancellation identity.}
From \eqref{eq:r-ode}:
$(\log r_n)' = G(x)\cos(2\theta_n) + D(x)$.
From Definition~\ref{def:wkb-amp}:
$(\log r_n^{\mathrm{WKB}})' = -(\wncl)'/(2\wncl)$.
Since $G = (\wncl)'/(2\wncl) + D$,
we have $-(\wncl)'/(2\wncl) = D - G$.
Subtracting (with
$\log(r_n/r_n^{\mathrm{WKB}})\big|_{x=x_0} = 0$):
\begin{equation}\label{eq:drift-cancel}
  \frac{d}{dx}\log\frac{r_n(x)}{r_n^{\mathrm{WKB}}(x)}
  = \bigl[G\cos(2\theta_n) + D\bigr]
    - \bigl[D - G\bigr]
  = G(x)\bigl(1 + \cos(2\theta_n(x))\bigr).
\end{equation}
The $D = p'/(4p)$ terms cancel algebraically as a
consequence of the definition of $r_n^{\mathrm{WKB}}$
as the solution of the averaged Pr\"ufer equation;
this is an exact algebraic identity, not a smallness
argument.

\textit{Step 2: Splitting.}
Write $G(1+\cos(2\theta_n)) = G + G\cos(2\theta_n)$.
Integrating \eqref{eq:drift-cancel} from $x_0$:
\begin{equation}\label{eq:log-ratio-split}
  \log\frac{r_n(x)}{r_n^{\mathrm{WKB}}(x)}
  = \int_{x_0}^x G(t)\,dt
    + \int_{x_0}^x G(t)\cos(2\theta_n(t))\,dt.
\end{equation}
Since $(\log r_n^{\mathrm{WKB}})' = D - G$,
integrating from $x_0$ to $x$ gives:
\[
  \int_{x_0}^x G(t)\,dt
  = \int_{x_0}^x D(t)\,dt
    - \log\frac{r_n^{\mathrm{WKB}}(x)}{r_n^{\mathrm{WKB}}(x_0)}.
\]
Substituting into \eqref{eq:log-ratio-split} and using
$r_n(x_0) = r_n^{\mathrm{WKB}}(x_0)$:
\begin{equation}\label{eq:drift-final}
  \log\frac{r_n(x)}{r_n^{\mathrm{WKB}}(x)}
  = \int_{x_0}^x G(t)\cos(2\theta_n(t))\,dt.
\end{equation}
Equation~\eqref{eq:drift-final} is the precise form of the
drift cancellation: the non-oscillatory mean term
$\int G\,dt$ cancels exactly against the WKB correction,
leaving only the oscillatory integral on the right-hand side.

\textit{Step 3: Exponentiation ---
proof of \eqref{eq:amp-ratio}.}
By Lemma~\ref{lem:prufer-oscillation-control}
(sup-form \eqref{eq:osc-control-sup} with $h = G$,
using $\|G\|_{L^\infty} \le C_{\delta,\gamma}/n$
and $\|G'\|_{L^1} \le C_{\delta,\gamma}/n$
from Lemma~\ref{lem:G-regularity}(iii),(iv)):
\begin{equation}\label{eq:sup-Gcos}
  \sup_{x \in I_\delta}
  \left|\int_{x_0}^x G(t)\cos(2\theta_n(t))\,dt\right|
  \le \frac{C_{\delta,\gamma}}{n}
      \left(\frac{C_{\delta,\gamma}}{n}
            + \frac{C_{\delta,\gamma}}{n}\right)
  = \frac{C_{\delta,\gamma}}{n}.
\end{equation}
From \eqref{eq:drift-final} and \eqref{eq:sup-Gcos}:
$\sup_{x \in I_\delta}
|\log(r_n(x)/r_n^{\mathrm{WKB}}(x))|
\le C_{\delta,\gamma}/n$.
Exponentiating:
$r_n(x)/r_n^{\mathrm{WKB}}(x) = 1 + O(1/n)$
uniformly on $I_\delta$,
and squaring gives \eqref{eq:amp-ratio}.

\textit{Step 4: Gronwall estimate, bootstrap closure,
and proof of \eqref{eq:amp-dev-deriv}.}
By Lemma~\ref{lem:deviation-ode}, $\Delta_n$ satisfies
\eqref{eq:deviation-ode} with
$|2D| \le C_\delta$,
$|\widetilde{B}_n| \le C_{\delta,\gamma}/n$,
$|G| \le C_{\delta,\gamma}/n$,
and $\Delta_n(x_0) = 0$.
The forcing term
$|2r_n^{\mathrm{WKB},2}\widetilde{B}_n|
\le C_{\delta,\gamma}r_n(x_0)^2/n$.
A Gronwall estimate on $I_\delta$
(length $\le 2T$, $\Delta_n(x_0) = 0$) yields:
\begin{equation}\label{eq:Delta-gronwall}
  |\Delta_n(x)| \le \frac{C_{\delta,\gamma}\,r_n(x_0)^2}{n}
  \qquad \text{for all } x \in I_\delta.
\end{equation}

We now verify bootstrap consistency:
the term $2\Delta_n G\cos(2\theta_n)$ in
\eqref{eq:deviation-ode} is self-referential,
since it involves $\Delta_n$ itself.
Applying Lemma~\ref{lem:prufer-oscillation-control}
to the integrated form of this term,
with $h = \Delta_n G/r_n(x_0)^2 \in C^1(I_\delta)$
and using \eqref{eq:Delta-gronwall} together with
Lemma~\ref{lem:G-regularity}(iii),(iv):
\[
  \|h\|_{L^\infty(I_\delta)}
  \le \frac{C_{\delta,\gamma}}{n}
      \cdot \frac{C_{\delta,\gamma}}{n}
  = \frac{C_{\delta,\gamma}}{n^2},
  \qquad
  \|h'\|_{L^1(I_\delta)} \le \frac{C_{\delta,\gamma}}{n^2}.
\]
Hence:
\begin{equation}\label{eq:bootstrap-check}
  \left|\int_{I_\delta}
    \Delta_n(t)\,G(t)\cos(2\theta_n(t))\,dt\right|
  \le \frac{C_{\delta,\gamma}\,r_n(x_0)^2}{n^2}
  = o\!\left(\frac{r_n(x_0)^2}{n}\right).
\end{equation}
This yields a contraction at the level of the error term:
the self-referential contribution \eqref{eq:bootstrap-check}
is strictly dominated by the main $O(r_n(x_0)^2/n)$ forcing.
In particular, no circular reasoning arises
in the Gronwall step \eqref{eq:Delta-gronwall}.
Substituting \eqref{eq:Delta-gronwall} back into
\eqref{eq:deviation-ode} and using
$|2\Delta_n G\cos(2\theta_n)|
\le C_{\delta,\gamma}r_n(x_0)^2/n^2$
establishes \eqref{eq:amp-dev-deriv}.
\end{proof}

\begin{remark}[Structure of the proof]
Steps~1--3 require no Gronwall argument:
the identity \eqref{eq:drift-final} reduces
$\log(r_n/r_n^{\mathrm{WKB}})$ to a single oscillatory
integral, controlled by one application of IBP.
Step~4 applies Gronwall only to $\Delta_n$,
whose forcing is $O(1/n)$ by
Lemma~\ref{lem:deviation-ode},
and closes the bootstrap via \eqref{eq:bootstrap-check},
which shows the self-referential term is $O(1/n^2)$,
strictly smaller than the main $O(1/n)$ bound.
\end{remark}

%%=================================================================
\section{Normalization Matching}\label{sec:normalization}
%%=================================================================

\begin{lemma}[WKB normalization]\label{lem:normalization}
For $n \le \gamma N$:
\begin{align}
  \frac{1}{2}\int_{I_\delta} r_n^{\mathrm{WKB}}(x)^2\,dx
  &= \lambda_n\left(1 + O\!\left(\frac{1}{n}\right)\right),
  \label{eq:norm-match}\\
  \frac{r_n^{\mathrm{WKB}}(x)^2}{2}
  &= \lambda_n\,\rhocl(x)
     \left(1 + O\!\left(\frac{1}{n}\right)\right).
  \label{eq:rn-vs-rhocl}
\end{align}
\end{lemma}

\begin{proof}
\textit{Step 1: Global normalization of $r_n^2$.}
Using $\sin^2\theta_n = \tfrac{1}{2} - \tfrac{1}{2}\cos(2\theta_n)$:
\[
  \psin(x)^2 = r_n(x)^2\sin^2\theta_n(x)
  = \frac{1}{2}r_n(x)^2
    - \frac{1}{2}r_n(x)^2\cos(2\theta_n(x)).
\]
Integrating over $I_\delta$ and using
Lemma~\ref{lem:turning-point-cutoff}:
\[
  \lambda_n
  = \int_{I_\delta}\psin^2 + O(e^{-\beta n}\lambda_n)
  = \frac{1}{2}\int_{I_\delta}r_n^2
    - \frac{1}{2}\int_{I_\delta}r_n^2\cos(2\theta_n)
    + O(e^{-\beta n}\lambda_n).
\]
For the oscillatory term: set $b_n = r_n^2/r_n(x_0)^2$.
By \eqref{eq:amp-ratio}, $\|b_n\|_{L^\infty} \le C$,
and by the Remark~\ref{rem:two-bounds} after
Lemma~\ref{lem:amplitude-drift},
$\|b_n'\|_{L^\infty} = O(1/n)$.
Applying Corollary~\ref{cor:osc-slowly-varying}
with $a \equiv 1$:
\[
  \left|\int_{I_\delta} r_n^2\cos(2\theta_n)\right|
  = r_n(x_0)^2
    \left|\int_{I_\delta} b_n\cos(2\theta_n)\right|
  \le \frac{C_{\delta,\gamma}\,r_n(x_0)^2}{n}.
\]
Since $r_n(x_0)^2 \le C\lambda_n$
(which follows from the $L^2$-normalization and
\eqref{eq:amp-ratio}):
\[
  \frac{1}{2}\int_{I_\delta} r_n^2
  = \lambda_n\bigl(1 + O(1/n)\bigr).
\]
Combined with \eqref{eq:amp-ratio}:
\[
  \frac{1}{2}\int_{I_\delta} r_n^{\mathrm{WKB}}(x)^2\,dx
  = \lambda_n\bigl(1 + O(1/n)\bigr),
\]
which is \eqref{eq:norm-match}.

\textit{Step 2: Pointwise ratio to $\rhocl$.}
From \eqref{eq:wkb-amp}:
\[
  r_n^{\mathrm{WKB}}(x)^2
  = r_n(x_0)^2\,\frac{\wncl(x_0)}{\wncl(x)}.
\]
Integrating over $I_\delta$ and using \eqref{eq:norm-match}:
\[
  r_n(x_0)^2\,\wncl(x_0)
  = \frac{2\lambda_n\bigl(1+O(1/n)\bigr)}
         {\displaystyle\int_{I_\delta} \frac{dt}{\wncl(t)}}.
\]

We now verify that $h(x) := 1/\wncl(x)$ satisfies the
regularity conditions of
Lemma~\ref{lem:prufer-oscillation-control}.
On $I_\delta$, $\wncl(x) \ge c_{\delta,\gamma}\,n$
by Lemma~\ref{lem:G-regularity}(i),
so:
\[
  \|h\|_{L^\infty(I_\delta)}
  \le \frac{1}{c_{\delta,\gamma}\,n},
  \qquad
  \|h'\|_{L^1(I_\delta)}
  = \left\|\frac{(\wncl)'}{\wncl^2}\right\|_{L^1(I_\delta)}
  \le \frac{C_{\delta,\gamma}}{n^2}\cdot |I_\delta|
  = O\!\left(\frac{1}{n^2}\right).
\]
This lower bound on $\wncl$ uses only that $x \in I_\delta$,
i.e., that the turning points $x_\pm(n)$ have been excluded
by the bulk cutoff of Lemma~\ref{lem:turning-point-cutoff};
the bound degenerates as $|x| \to x_+(n)$,
but this region has already been removed.

Hence $r_n^{\mathrm{WKB}}(x)^2/2$ equals:
\[
  \frac{\lambda_n\bigl(1+O(1/n)\bigr)}
       {\wncl(x)\cdot
        \displaystyle\int_{I_\delta} \frac{dt}{\wncl(t)}}
  = \lambda_n\,\rhocl(x)\bigl(1+O(1/n)\bigr),
\]
using Lemma~\ref{lem:bohr-sommerfeld-pswf}
and Definition~\ref{def:rho-cl},
which gives \eqref{eq:rn-vs-rhocl}.
\end{proof}

%%=================================================================
\section{Proof of the Main Theorem}\label{sec:weaklimit}
%%=================================================================

\begin{theorem}[Semiclassical equidistribution of PSWF densities]%
\label{thm:weak-limit}
For $n \le \gamma N$ and every $f \in C^1([-T,T])$:
\begin{equation}\label{eq:weak-limit}
  \int_{-T}^T f(x)\,\psin(x)^2\,dx
  = \lambda_n(c)
    \int_{-x_+(n)}^{x_+(n)} f(x)\,\rhocl(x)\,dx
    + R_n(f),
\end{equation}
where:
\begin{equation}\label{eq:remainder}
  |R_n(f)|
  \le \frac{C_{\delta,\gamma}\,\lambda_n(c)}{n}
      \left(\|f\|_{L^\infty} + T\,\|f'\|_{L^\infty}\right).
\end{equation}
The constant $C_{\delta,\gamma}$ depends only on
$\delta$, $\gamma$, $c$, $T$, not on $n$ or $f$.
\end{theorem}

\begin{proof}
\textit{Step 1 (Bulk cutoff).}
By Lemma~\ref{lem:turning-point-cutoff}:
\[
  \int_{-T}^T f\psin^2
  = \int_{I_\delta} f\psin^2
    + O\!\left(e^{-\beta n}\lambda_n\|f\|_{L^\infty}\right).
\]

\textit{Step 2 (Pr\"ufer decomposition).}
Using $\sin^2\theta_n = \frac{1}{2} - \frac{1}{2}\cos(2\theta_n)$:
\[
  \int_{I_\delta} f\psin^2
  = \frac{1}{2}\int_{I_\delta} f\,r_n^2
    - \frac{1}{2}\int_{I_\delta} f\,r_n^2\cos(2\theta_n).
\]

\textit{Step 3 (Oscillation term).}
Set $a = f \in C^1(I_\delta)$
and $b_n = r_n^2/r_n(x_0)^2$.
By Remark~\ref{rem:two-bounds}:
$\|b_n\|_{L^\infty} \le C$ and
$\|b_n'\|_{L^\infty} = O(1/n)$.
Applying Corollary~\ref{cor:osc-slowly-varying}:
\[
  \left|\int_{I_\delta} f\,r_n^2\cos(2\theta_n)\right|
  \le \frac{C_{\delta,\gamma}\,\lambda_n}{n}
      \bigl(\|f\|_{L^\infty} + T\|f'\|_{L^\infty}\bigr).
\]

\textit{Step 4 (Mean term).}
By \eqref{eq:amp-ratio} and \eqref{eq:rn-vs-rhocl}:
\[
  r_n(x)^2 = r_n^{\mathrm{WKB}}(x)^2\bigl(1+O(1/n)\bigr)
  = 2\lambda_n\,\rhocl(x)\bigl(1+O(1/n)\bigr).
\]
Hence:
\[
  \frac{1}{2}\int_{I_\delta} f\,r_n^2
  = \lambda_n\int_{-x_+(n)}^{x_+(n)} f\,\rhocl
    + O\!\left(\frac{\lambda_n\|f\|_{L^\infty}}{n}\right).
\]

\textit{Step 5 (Boundary terms).}
The contribution on $[-T,T]\setminus I_\delta$ is
$O(e^{-\beta_\gamma n}\lambda_n\|f\|_{L^\infty})$
by Lemma~\ref{lem:turning-point-cutoff}.
The integral
$\int_{I_\delta\setminus[-x_+(n),x_+(n)]}f\,\rhocl = 0$
since $\rhocl$ is supported on
$[-x_+(n),x_+(n)] \subset I_\delta$
(Definition~\ref{def:rho-cl},
valid since $x_+(n) \le (1-\eta_\gamma)T < (1-\delta)T$
for $\delta < \eta_\gamma$).

\textit{Assembly.}
Combining Steps~1--5:
\begin{align*}
  \int_{-T}^T f\psin^2
  &= \lambda_n\int_{-x_+(n)}^{x_+(n)} f\,\rhocl
     + O\!\left(\frac{\lambda_n}{n}
       \bigl(\|f\|_{L^\infty}+T\|f'\|_{L^\infty}\bigr)\right)
     + O\!\left(e^{-\beta n}\lambda_n\|f\|_{L^\infty}\right).
\end{align*}
Since $e^{-\beta n} = O(1/n)$,
this gives \eqref{eq:weak-limit}--\eqref{eq:remainder}
with $C_{\delta,\gamma}$ independent of $n$ and $f$.
\end{proof}

\begin{remark}[Uniformity of the constant]
The constant satisfies
$C_{\delta,\gamma} \lesssim c_{\delta,\gamma}^{-2}$,
where $c_{\delta,\gamma}$ is the phase lower bound
constant from Lemma~\ref{lem:phase-lower}.
It is uniform in $n \le \gamma N$ and in $f \in C^1$
for fixed $\delta$, $\gamma$, $c$, $T$.
\end{remark}

\begin{corollary}[Closure of the bulk program]%
\label{cor:bulk-program-closed}
Theorem~\ref{thm:weak-limit} supplies hypothesis~(6.1) of
Lemma~6.3 of \cite{PaperIII} with
$\omega_n = C_{\delta,\gamma}/n$,
completing the logical chain:
\[
  \text{Thm.~\ref{thm:weak-limit}}
  \;\Rightarrow\; \texttt{lem:bulk-reduction}
  \;\Rightarrow\; \texttt{prob:comm-refined}
  \;\Rightarrow\; \text{Bulk tail bound}
\]
for all $n \le \gamma N$,
conditional on Assumption~2.4 of \cite{PaperIII}.
\end{corollary}

%%=================================================================
\section{Open Problems}\label{sec:open}
%%=================================================================

\begin{enumerate}[(P1)]
  \item \textbf{Extension to $f \in C^0$.}
  Theorem~\ref{thm:weak-limit} requires $f \in C^1$.
  For $f \in C^0$, a density argument gives weak-$*$
  convergence without a rate.
  A rate for $f \in C^\alpha$ ($0 < \alpha < 1$)
  via interpolation between $C^0$ and $C^1$ would extend
  the result to H\"older-continuous test functions.

  \item \textbf{Off-diagonal regime.}
  An analogue of Theorem~\ref{thm:weak-limit} for
  $|m-n| \ge \delta N$ (cross terms $\psi_m\psi_n$)
  remains open and would complete the off-diagonal
  program of \cite{PaperIII}.

  \item \textbf{Explicit Bohr--Sommerfeld constant.}
  Making the implicit constant in \eqref{eq:BS}
  fully explicit in terms of $c$, $T$, $\gamma$
  would allow the remainder bound in
  Theorem~\ref{thm:weak-limit} to be computed
  numerically for given parameters.

  \item \textbf{Uniformity in $c$.}
  All results treat $c$ and $T$ as fixed.
  The regime $c \to \infty$ with $n/c$ fixed
  requires uniform WKB estimates in $c$
  that are not directly available from
  \cite{OsipovRokhlinXiao} as cited here.

  \item \textbf{Near-critical regime $n \sim \gamma N$.}
  As $n \to \gamma N$, the bulk interval $I_\delta$
  shrinks to a point and the present argument degenerates.
  A separate analysis, likely requiring Airy-function
  scaling near the spectral edge,
  is needed to handle this regime.
\end{enumerate}

%%=================================================================
\begin{thebibliography}{99}

\bibitem{PaperI}
Anonymous Author,
\emph{Coercivity of the Prolate Gram Form at Discrete
Sampling Points},
preprint, 2026.

\bibitem{PaperII}
Anonymous Author,
\emph{Scaling Limits of Prolate Gram Forms, Uniform
Coercivity, and a Trace Formula Toward Weil Positivity},
preprint, 2026.

\bibitem{PaperIII}
Anonymous Author,
\emph{PSWF Product Spectral Tail Estimates: Bulk and
Off-Diagonal Regimes, and an Obstruction at the
Spectral Edge},
preprint, 2026.

\bibitem{Levitan}
B.~M.~Levitan and I.~S.~Sargsjan,
\textit{Sturm--Liouville and Dirac Operators},
Kluwer, Dordrecht, 1991.

\bibitem{OsipovRokhlinXiao}
A.~Osipov, V.~Rokhlin, and H.~Xiao,
\textit{Prolate Spheroidal Wave Functions of Order Zero},
Springer, New York, 2013.

\bibitem{Slepian1978}
D.~Slepian,
\textit{Prolate spheroidal wave functions, Fourier
analysis, and uncertainty~-- V},
Bell Syst.\ Tech.\ J.\ \textbf{57} (1978), 1371--1430.

\bibitem{Szego}
G.~Szeg\H{o},
\textit{Orthogonal Polynomials},
AMS Colloquium Publications, vol.~23, 1939.

\bibitem{Zygmund}
A.~Zygmund,
\textit{Trigonometric Series},
Cambridge University Press, 3rd~ed., 2002.

\end{thebibliography}

\end{document}
